% 学术学位硕士研究生中期考核空白模板。
% 复制本文件后，先修改下列基本信息，再填写各报告部分。
\documentclass[ma]{xdupgtm}
% 学号
\newcommand\studentNo{}
% 论文题目，可使用两行
\newcommand\thesisTitleOne{}
\newcommand\thesisTitleTwo{}
% 姓名、一级学科、二级学科、指导教师和学院
\newcommand\authorName{}
\newcommand\major{}
\newcommand\submajor{}
\newcommand\advisorName{}
\newcommand\school{}
% 报告日期和选题来源
\newcommand\submitDate{}
\newcommand\subjectSource{}
% 需要电子签名时取消注释，并将图片放入 figures/sign/。
% \renewcommand\zdjsqm{\sign{advisor.png}}
% \renewcommand\zdjsqmrq{年\qquad 月\qquad 日}
% \renewcommand\zzqm{\sign{chair.png}}
% \renewcommand\cyqm{\sign{member-one.png}\emptysign\sign{member-two.png}}
% \renewcommand\cyqmrq{年\qquad 月\qquad 日}

\begin{document}
% 章节命令仅供编辑器大纲使用，条件分支不会输出到 PDF。
\iffalse
\section{选题来源}
\fi
% 选题来源只需在 xzlyj 环境外修改 subjectSource。
\begin{xzlyj}\end{xzlyj}
\iffalse
\section{学位论文研究目标及研究内容}
\fi
\begin{yjmbnr}
% 在这里填写学位论文研究目标及研究内容；应与开题报告保持一致。
\end{yjmbnr}
\iffalse
\section{目前已完成学位论文工作的内容}
\fi
\begin{ywcgz}
% 在这里填写目前已完成的理论、实验研究内容及获得的结论。
\end{ywcgz}
\iffalse
\section{中期考核工作计划}
\fi
\begin{workplan}
% 第三至第五部分合用一个表格，保留 workplan 外层环境。
\iffalse
\subsection{现阶段完成的工作与开题报告内容不相符的情况说明}
\fi
\begin{bxfsm}
% 在这里填写与开题报告不相符的情况；没有变化时填写“无”。
\end{bxfsm}
\iffalse
\subsection{下一步工作计划及需要完成的研究内容和需要解决的关键技术}
\fi
\begin{gzhjh}
% 在这里填写下一步工作计划、拟完成的研究内容和需要解决的关键技术。
\end{gzhjh}
\iffalse
\subsection{已发表的与学位论文相关的学术论文、其他研究成果以及拟发表的研究成果}
\fi
\begin{kycg}
% 在这里填写已发表、已录用或拟发表的论文及其他研究成果。
\end{kycg}
\end{workplan}
\iffalse
\section{指导教师评价}
\fi
\begin{jsyj}
% 在这里填写指导教师评价意见。
\end{jsyj}
\iffalse
\section{中期总结报告评语及结论}
\fi
\begin{pyjl}
% 在这里填写中期总结报告评语及“通过”或“不通过”结论。
\end{pyjl}
\end{document}
